\documentclass{../../../templates/natsirt}

\usepackage{pgfplots}
\pgfplotsset{compat=1.18}

\title{Ordinary Differential Equations}
\author{Yi Hong}
\subtitle{Note:p}

\setlength{\parindent}{0pt}
\setlength{\parskip}{3pt}
\setlength{\headheight}{15pt}

% Custom Box Environments for Pure Active-Recall Checklist
\newcounter{topic}[section]

\newenvironment{recallbox}[1][]{%
    \begin{bluebox}[#1]\refstepcounter{topic}\textbf{Definitions \& Concepts to State \thesection.\thetopic}%
}{\end{bluebox}}

\newenvironment{theoremsbox}[1][]{%
    \begin{greenbox}[#1]\refstepcounter{topic}\textbf{Theorems \& Identities to State \thesection.\thetopic}%
}{\end{greenbox}}

\newenvironment{challengebox}[1][]{%
    \begin{redbox}[#1]\refstepcounter{topic}\textbf{Proof Challenges \& Calculations \thesection.\thetopic}%
}{\end{redbox}}

\newenvironment{warningbox}[1][]{\begin{whitebox}[#1]\textbf{Conventions \& Exam Pitfalls. }}{\end{whitebox}}
\newenvironment{reviewbox}[1][]{\begin{whitebox}[#1]\textbf{Textbook Exercise Checklist. }}{\end{whitebox}}

\newcommand{\citem}{\item[$\square$]}
\newcommand{\citemdone}{\item[$\text{\rlap{$\checkmark$}}\square$]}

\begin{document}

\section{Chapter 1: Introduction to ODEs}

\begin{recallbox}
\begin{itemize}
    \citem What is a direction field?
    \citem Write the definition of ODE vs.\ PDE (single vs.\ several independent variables), and what is a system of differential equations (e.g., Lotka--Volterra predator--prey equations).
    \citem Write the definition of the order of an ODE.
    \citem Write the definition of the normal form of an ODE.
    \begin{itemize}
        \item Remark: Not every differential equation can be written in normal form (e.g., $\sin(y'') + xy'' = y$ cannot be explicitly solved for $y''$).
    \end{itemize}
    \citem Write the definition of when an ODE is called linear.
    \citem Write the definition of a solution of an ODE.
\end{itemize}
\end{recallbox}

\begin{challengebox}
\begin{itemize}
    \citem Model a falling object using Newton's second law ($m\frac{dv}{dt} = mg - \gamma v$), find the equilibrium solution ($\frac{dv}{dt} = 0$), and explain what terminal velocity is.
    \citem Verify a solution defined by a definite integral (using the Fundamental Theorem of Calculus): verify that $y(t) = e^{t^2}\int_0^t e^{-s^2}\,ds + e^{t^2}$ satisfies $y' - 2ty = 1$.
\end{itemize}
\end{challengebox}

\section{Chapter 2: First-Order Linear ODEs and Solution Behavior}

\begin{warningbox}
\begin{itemize}
    \citem Remember to rewrite $P(t)\frac{dy}{dt} + Q(t)y = G(t)$ into standard form $\frac{dy}{dt} + p(t)y = g(t)$ by dividing by $P(t) \neq 0$ before finding the integrating factor $\mu(t) = \exp\!\left(\int p(t)\,dt\right)$.
\end{itemize}
\end{warningbox}

\begin{recallbox}
\begin{itemize}
    \citem Write out the procedure of the method of integrating factor, and give the explicit general solution formula $y = \frac{1}{\mu(t)}\!\left(\int_{t_0}^t \mu(s)g(s)\,ds + c\right)$.
    \citem When a question asks to analyze the behavior of a solution, analyze the properties of $y'$: when it is positive, when it is negative, what its zeros are, and when $y$ or $y'$ will diverge or converge.
    \citem For analyzing the behavior of a solution, another direct approach is to simply solve the equation.
\end{itemize}
\end{recallbox}

\begin{challengebox}
\begin{itemize}
    \citem How do we solve a first-order linear ODE? For example, solve
    \[
        (4 + t^2)\frac{dy}{dt} + 2ty = 4t.
    \]
\end{itemize}
\end{challengebox}

\section{Lecture 2: Separable, Autonomous, and Exact Equations}

\begin{warningbox}
\begin{itemize}
    \citem Sometimes using an integrating factor to solve a question is better than separating the equation.
    \citem If you use the separable equation method to solve an ODE, always check whether there are singular solutions.
\end{itemize}
\end{warningbox}

\begin{recallbox}
\begin{itemize}
    \citem What is an initial value problem (IVP)?
    \citem What is a separable equation?
    \citem Write the definition of an autonomous differential equation.
    \citem What is an equilibrium solution of an autonomous equation, and what does it mean for an equilibrium solution $y_0$ of $y' = f(y)$ to be:
    \begin{itemize}
        \item stable (if $f(y)$ changes sign from $+$ to $-$ at $y_0$);
        \item unstable (if $f(y)$ changes sign from $-$ to $+$ at $y_0$);
        \item semistable (if $f(y)$ does not change sign at $y_0$)?
    \end{itemize}
    \citem State the procedure for drawing a phase line portrait for an autonomous equation $y' = f(y)$.
    \citem Write the definition of an exact differential equation, along with the exactness conditions.
    \citem How do we make a non-exact equation exact using an integrating factor $\mu(x,y)$?
\end{itemize}
\end{recallbox}

\begin{challengebox}
\begin{itemize}
    \citem Refer to the modeling examples mentioned in the lecture notes: the mixing problem, population growth with predation (field mice and owls, $\frac{dp}{dt} = \frac{p}{2} - 450$), and escape velocity / variable gravitational field ($m\frac{dv}{dt} = -\frac{mgR^2}{(R+x)^2}$, using $v\frac{dv}{dx}$ to find $v(x)$, maximum altitude $A_{\max}$, and escape velocity $\sqrt{2gR}$).
    \citem Solve the differential equation
    \[
        (3x^2 - 2xy + 2)\,dx + (6y^2 - x^2 + 3)\,dy = 0.
    \]
    \citem Check whether $x^2y^3 + x(1 + y^2)y' = 0$ is exact, then multiply by $\mu(x,y) = \frac{1}{xy^3}$ and solve it.
\end{itemize}
\end{challengebox}

\section{Lecture 3: Existence and Uniqueness, and Second-Order Homogeneous ODEs}

\begin{recallbox}
\begin{itemize}
    \citem Write the definition of a homogeneous differential equation and a nonhomogeneous differential equation.
    \citem Write the definition of the associated homogeneous equation.
    \citem Write the definition of the characteristic equation (or auxiliary equation), and explain how to use it to solve a second-order differential equation.
    \citem For the differential operator of a second-order differential equation, note that it is a linear operator; write the formal definition of a linear operator ($L[y_1 + y_2] = L[y_1] + L[y_2]$ and $L[cy] = cL[y]$) and symbolic differential operator notation ($Dy = y', \; D^2y = y'', \dots$).
    \citem Write the definition of a fundamental solution set.
    \citem How do we solve a second-order homogeneous DE when the discriminant of the characteristic equation is negative (using Euler's formula $e^{a+ib} = e^a(\cos b + i\sin b)$ and taking the real and imaginary parts of $e^{(\alpha + i\beta)t}$ to get $\{e^{\alpha t}\cos(\beta t), \, e^{\alpha t}\sin(\beta t)\}$)?
    \citem If a homogeneous second-order DE has a repeated root ($r = -\frac{b}{2a}$) for the characteristic equation, how do we use the reduction of order method ($y(t) = v(t)y_1(t)$) to find the second solution and build the fundamental solution set?
\end{itemize}
\end{recallbox}

\begin{theoremsbox}
\begin{itemize}
    \citem First-order Existence and Uniqueness Theorem: If $f$ and $\frac{\partial f}{\partial y}$ are continuous in a rectangle $R : |t| \le a, \; |y| \le b$, then there exists some interval $|t| \le h \le a$ in which there exists a unique solution
    \[
        y = \phi(t)
    \]
    of the initial value problem $y' = f(t, y), \; y(0) = 0$.
    \begin{itemize}
        \item Remark: Continuity of $f$ alone guarantees existence of a solution in some open interval, while continuity of $\frac{\partial f}{\partial y}$ is needed to guarantee uniqueness.
        \item Write the equivalent integral equation $\phi(t) = \int_0^t f(s, \phi(s))\,ds$ and the general Picard iteration formula $\phi_{n+1}(t) = \int_0^t f(s, \phi_n(s))\,ds$ with $\phi_0(t) = 0$; explain why the interval is restricted to $h = \min(a, b/M)$ where $|f(t, y)| \le M$ on $R$.
    \end{itemize}
    \citem Write the Principle of Superposition, and prove it.
    \citem Write the Existence and Uniqueness Theorem for second-order differential equations.
    \citem Write the definition of the Wronskian
    \[
        W[y_1, y_2](t) = \begin{vmatrix} y_1(t) & y_2(t) \\ y_1'(t) & y_2'(t) \end{vmatrix} = y_1(t)y_2'(t) - y_1'(t)y_2(t)
    \]
    and the theorem that $W[y_1, y_2](t_0) \neq 0$ implies every solution can be written as $y(t) = C_1y_1(t) + C_2y_2(t)$.
    \citem Write Abel's Theorem, and prove it.
\end{itemize}
\end{theoremsbox}

\begin{challengebox}
\begin{itemize}
    \citem Use Picard's iteration method to find $\phi_1(t)$, $\phi_2(t)$, and $\phi_3(t)$ for the IVP $y' = 2t(1 + y), \; y(0) = 0$.
    \citem Determine whether the Existence and Uniqueness Theorem guarantees a unique solution for $y' = 3y^{2/3}$ with:
    \begin{itemize}
        \item (a) $y(2) = 0$
        \item (b) $y(2) = 5$
    \end{itemize}
    \citem Find the largest interval of existence and uniqueness for the IVP
    \[
        e^t y'' - \frac{y'}{t-3} + y = \ln t, \quad y(1) = y_0, \quad y'(1) = y_1.
    \]
    \citem If $W[f, g](t) = 3e^{4t}$ and $f(t) = e^{2t}$, find $g(t)$.
    \citem Solve $y'' - 2y' + 2y = 0$.
\end{itemize}
\end{challengebox}

\section{Lecture 4: Second-Order Nonhomogeneous ODEs (Undetermined Coefficients)}

\begin{warningbox}
\begin{itemize}
    \citem Remark: For a second-order ODE with real coefficients, complex conjugate roots $\alpha \pm i\beta$ can only be simple ($s = 1$); repeated complex roots require order $\ge 4$.
\end{itemize}
\end{warningbox}

\begin{recallbox}
\begin{itemize}
    \citem Write the procedure for solving a second-order nonhomogeneous ODE (with real coefficients).
    \citem Suppose for the ODE
    \[
        ay'' + by' + cy = g(t)
    \]
    we have 4 cases for $g(t)$; write down these four cases and the corresponding form of the particular solution (including the multiplicity factor $t^s$, where $s = 0, 1, 2$ is the multiplicity of the root in the characteristic equation, and $m = \max\{m_1, m_2\}$ with both $\cos(\beta t)$ and $\sin(\beta t)$ terms).
    \citem How do we solve when $g(t)$ is a sum of terms, such as $g(t) = e^t + \cos 3t$?
\end{itemize}
\end{recallbox}

\begin{challengebox}
\begin{itemize}
    \citem Solve:
    \begin{itemize}
        \item $y'' + y' = 1 - 2t^2$
        \item $y'' - 4y' + 4y = 16e^{2t}$
        \item $y'' - 3y' + 2y = 4t\sin t$
        \item $y'' - 9y = e^{3t}\cos t$
    \end{itemize}
    \citem Solve
    \[
        y'' - 3y' - 4y = 3e^{2t} + 2\sin t - 8e^t\cos(2t).
    \]
\end{itemize}
\end{challengebox}

\section{Week 5: Variation of Parameters and Simple Spring-Mass Systems}

\begin{warningbox}
\begin{itemize}
    \citem Remember to put the ODE into standard form $y'' + p(t)y' + q(t)y = g(t)$ first before applying variation of parameters.
\end{itemize}
\end{warningbox}

\begin{recallbox}
\begin{itemize}
    \citem Write the procedure for using the variation of parameters method to solve a second-order ODE, and prove the formula (impose the auxiliary condition $u_1'(t)y_1(t) + u_2'(t)y_2(t) = 0$, and give the explicit formulas $u_1(t) = -\int_{t_0}^t \frac{g(s)y_2(s)}{W(y_1, y_2)(s)}\,ds$ and $u_2(t) = \int_{t_0}^t \frac{g(s)y_1(s)}{W(y_1, y_2)(s)}\,ds$).
    \citem Remark on the advantages and difficulties of variation of parameters: it works for arbitrary $g(t)$ and variable coefficients once $\{y_1, y_2\}$ is known; the difficulties are finding $\{y_1, y_2\}$ for non-constant coefficients and evaluating the integrals.
    \citem Model a simple spring-mass system (using equilibrium balance $mg - kL = 0$, i.e., $k = mg/L$, and the four forces $F_g, F_s, F_d, F_{\text{ex}}$ leading to $mu'' + \gamma u' + ku = F_{\text{ex}}(t)$); then ignore the damping force and external force and derive the amplitude ($R = \sqrt{c_1^2 + c_2^2}$), phase angle ($\tan\delta = c_2/c_1$ in $u(t) = R\cos(\omega_0 t - \delta)$), natural frequency, and period ($T = 2\pi\sqrt{m/k}$, and how $m$ and $k$ affect $T$) of the system.
\end{itemize}
\end{recallbox}

\begin{challengebox}
\begin{itemize}
    \citem Find the solution of
    \[
        y'' + 4y = 8\tan t.
    \]
\end{itemize}
\end{challengebox}

\section{Week 6: Damped and Forced Vibrations}

\begin{warningbox}
\begin{itemize}
    \citem Remark on all three cases ($\gamma^2 - 4km > 0, \; = 0, \; < 0$) for damped free vibrations: in every case $u(t) \to 0$ as $t \to \infty$; as $\gamma \to 2\sqrt{km}$, $\mu \to 0$ and $T_d \to \infty$; for critically damped ($\gamma = 2\sqrt{km}$) and overdamped ($\gamma > 2\sqrt{km}$) motion, the mass passes through equilibrium at most once.
\end{itemize}
\end{warningbox}

\begin{recallbox}
\begin{itemize}
    \citem For a vibration system, write the definition of quasi-frequency and quasi-period.
    \citem Compare the quasi-frequency with the frequency in undamped motion, explain whether small damping reduces or increases the frequency of oscillation, and compare the periods as well.
    \citem Write the definition of when a homogeneous spring-mass system is called critically damped and overdamped.
    \citem For a forced damped spring-mass system, write the form of the external force.
    \citem Write the definition of the transient solution for the forced damped spring-mass system.
    \citem Write the definition of the steady-state solution for the forced damped spring-mass system.
    \citem For forced damped vibrations $mu'' + \gamma u' + ku = F_0\cos(\omega t)$, write the steady-state solution $y_p(t) = R\cos(\omega t - \delta)$ with $R = \frac{F_0}{\Delta}$ and $\Delta = \sqrt{m^2(\omega_0^2 - \omega^2)^2 + \gamma^2\omega^2}$:
    \begin{itemize}
        \item check the limits of $R$ as $\omega \to 0$ ($R \to F_0/k$) and $\omega \to \infty$ ($R \to 0$);
        \item state the frequency of maximum amplitude $\omega_{\max}^2 = \omega_0^2\!\left(1 - \frac{\gamma^2}{2mk}\right)$, maximum amplitude $R_{\max} = \frac{F_0}{\gamma\omega_0\sqrt{1 - \gamma^2/(4mk)}}$, and resonance when $\gamma$ is small and $\omega \approx \omega_0$.
    \end{itemize}
    \citem Electric Circuits (Series RLC Circuit):
    \begin{itemize}
        \item state Kirchhoff's second law and the voltage drops across a resistor ($IR$), capacitor ($Q/C$), and inductor ($L\frac{dI}{dt}$), with $I = \frac{dQ}{dt}$;
        \item write the differential equation and initial conditions for charge $Q(t)$: $LQ'' + RQ' + \frac{1}{C}Q = E(t), \; Q(t_0) = Q_0, \; Q'(t_0) = I_0$;
        \item write the differential equation and initial conditions for current $I(t)$: $LI'' + RI' + \frac{1}{C}I = E'(t), \; I(t_0) = I_0, \; I'(t_0) = \frac{E(t_0) - RI_0 - Q_0/C}{L}$.
    \end{itemize}
    \citem Forced Vibrations Without Damping ($\gamma = 0$: Beats and Pure Resonance):
    \begin{itemize}
        \item for $mu'' + ku = F_0\cos(\omega t)$ with $\omega \neq \omega_0$ and $u(0) = u'(0) = 0$, derive $u(t) = \frac{2F_0}{m(\omega_0^2 - \omega^2)}\sin\!\left(\frac{(\omega_0 - \omega)t}{2}\right)\sin\!\left(\frac{(\omega_0 + \omega)t}{2}\right)$, state the amplitude envelope, and explain what a beat is when $|\omega_0 - \omega|$ is small;
        \item for $\omega = \omega_0$ (pure resonance), explain why $u_p(t) = t(A\cos\omega_0 t + B\sin\omega_0 t)$ has unboundedly growing amplitude as $t \to \infty$.
    \end{itemize}
\end{itemize}
\end{recallbox}

\begin{challengebox}
\begin{itemize}
    \citem A mass weighing 4 lb stretches a spring 2 in. Suppose that the mass is given an additional 6 in displacement in the positive direction and then released. The mass is in a medium that exerts a viscous resistance of 6 lb when the mass has a velocity of 3 ft/s. Under the assumptions discussed in this section, formulate the initial value problem that governs the motion of the mass.
    \citem Solve the damped vibration equation in terms of the parameters:
    \[
        mu'' + \gamma u' + ku = 0.
    \]
    \citem Solve $u'' + u = 0.5\cos(0.8t), \; u(0) = 0, \; u'(0) = 0$.
    \citem Solve $u'' + u = 0.5\cos t, \; u(0) = 0, \; u'(0) = 0$.
\end{itemize}
\end{challengebox}

\section{Week 7: Higher-Order Linear ODEs}

\begin{recallbox}
\begin{itemize}
    \citem Write the standard form of an $n$th-order linear ODE ($L[y] = y^{(n)} + p_1(t)y^{(n-1)} + \dots + p_n(t)y = g(t)$, dividing by $P_0(t) \neq 0$) and the $n$ initial conditions $y(t_0) = y_0, \dots, y^{(n-1)}(t_0) = y_0^{(n-1)}$.
    \citem Write the definition of the $n \times n$ Wronskian $W(y_1, \dots, y_n)(t)$, and the definition of linearly dependent and linearly independent functions $f_1, \dots, f_n$ on an interval $I$.
    \citem Explain why the general solution of the nonhomogeneous $n$th-order linear ODE is $y = c_1y_1(t) + \dots + c_ny_n(t) + Y(t)$ (using $L[Y_1 - Y_2] = 0$).
    \citem How do we solve higher-order linear ODEs with constant coefficients by undetermined coefficients? For the homogeneous part, there are three cases; state the procedure for dealing with these three cases.
    \citem Describe the asymptotic behavior of solutions to homogeneous constant-coefficient ODEs with distinct real roots as $t \to \infty$ (dominated by the algebraically largest root: positive $\to \pm\infty$, negative $\to 0$, zero $\to$ nonzero constant).
\end{itemize}
\end{recallbox}

\begin{theoremsbox}
\begin{itemize}
    \citem How do we determine the interval on which a solution exists for a higher-order linear ODE (Existence and Uniqueness Theorem for $n$th-order linear ODEs)?
    \citem Prove that if $W(y_1, \dots, y_n)(t_0) \neq 0$, then $y_1, \dots, y_n$ are linearly independent; state the theorem that solutions $y_1, \dots, y_n$ form a fundamental set of solutions on $I$ if and only if they are linearly independent on $I$.
\end{itemize}
\end{theoremsbox}

\begin{challengebox}
\begin{itemize}
    \citem Find the solution of
    \[
        y^{(4)} + 2y'' + y = 2\sin t.
    \]
    \citem Solve these higher-order ODE problems:
    \begin{itemize}
        \item $y^{(4)} - y = 0, \; y(0) = \frac{7}{2}, \; y'(0) = -4, \; y''(0) = \frac{5}{2}, \; y'''(0) = -2$
        \item $y^{(4)} + y = 0$
        \item $y''' - 3y'' + 3y' - y = 4e^t$
        \item $y''' - 4y' = t + 3\cos t + e^{-2t}$
    \end{itemize}
\end{itemize}
\end{challengebox}

\end{document}
